\section{\texorpdfstring{Softmax 交叉熵为何化简为 \(p-y\)}{Softmax 交叉熵为何化简为 p-y}}
\label{sec:10-Matrix-Calculus/04-Applications/02}

多分类模型对每个样本输出 logits \(z\in\R^K\)。softmax 把它变成概率向量：

\[
p_i = \frac{e^{z_i}}{\sum_j e^{z_j}}.
\]

正确类别用 one-hot \(y\) 标记（\(y_c=1\)，其余为 0），交叉熵损失：

\[
\ell(z,y) = -\sum_i y_i\log p_i.
\]

训练需要 \(\nabla_z\ell\)。初看之下，这是两个操作的复合：softmax 非线性（\(z\to p\)），再套交叉熵（\(p\to\ell\)）。链式法则说先算 softmax 的 Jacobian \(\partial p/\partial z\)（一个 \(K\times K\) 矩阵），再乘 \(\partial\ell/\partial p\)（一个 \(K\) 维向量），得到一个 \(K\) 维梯度。ImageNet \(K=1000\)，每步都在构造百万级矩阵去拿千维答案。这不是数学错，是计算浪费。

\subsection{走一趟笨路，看清中间发生什么}

动手走一遍。softmax 的偏导数（\(\delta_{ij}\) 是 Kronecker \(\delta\)）：

\[
\frac{\partial p_i}{\partial z_j} = p_i(\delta_{ij}-p_j).
\]

矩阵形式：

\[
J_{\mathrm{softmax}} = \diag(p) - pp^{\trans}.
\]

这是一个对角矩阵减去秩一外积——形式已算简洁，但仍然是 \(K^2\) 个元素。

交叉熵对概率的导数：\(\partial\ell/\partial p_i = -y_i/p_i\)。乘起来，逐分量：

\[
\begin{aligned}
\frac{\partial\ell}{\partial z_j}
&= \sum_i \frac{\partial\ell}{\partial p_i}\frac{\partial p_i}{\partial z_j} \\
&= \sum_i p_i(\delta_{ij}-p_j)\left(-\frac{y_i}{p_i}\right) \\
&= -\sum_i \delta_{ij}y_i + p_j\sum_i y_i \\
&= -y_j + p_j.
\end{aligned}
\]

最后一步用了 \(\sum_i y_i = 1\)（one-hot 标签分量之和为 1）。结果极度简洁：\textbf{梯度就是模型预测的概率减真实标签。}

这条路走通了，路径上两处隐患也一并暴露。第一，中间出现了 \(y_i/p_i\)。训练初期，真实类别的预测概率可能很小（比如均匀分布时 \(p_c\approx 1/K\)），\(y_c/p_c = K\)，不算问题。但一旦某个非目标类别的概率 \(p_i\) 趋近于零，对应的 \(y_i=0\)，\(0/0\) 型的未定义没有出现——真正的危险来自目标类别的 \(p_c\) 趋近于零时 \(y_c/p_c = 1/p_c\) 发散，幸好反向传播中我们从不单独计算 \(y_i/p_i\)，而是整体乘上 \(p_i(\delta_{ij}-p_j)\)，发散被抵消。但这个抵消仅在使用完整复合梯度时才可靠——如果你在计算图中把 softmax 和交叉熵拆成两个独立算子、分别计算局部梯度再乘，中间值可能已经在某一步越界。

第二处隐患：你构造了完整的 softmax Jacobian。即使在 \(\diag(p)-pp^{\trans}\) 这个简洁形式下，它仍然是 \(O(K^2)\) 存储。梯度公式只有一个对角的 \(p_j\) 减 \(y_j\)——凭什么绕那么远？把 softmax 和交叉熵拆成两层各自独立的链式乘法，是在不该拆的地方拆了一刀。

\subsection{换条路：把损失写成 logits 的直接函数}

交叉熵的结构暗示了更短的路。写出 \(\log p_i\)：

\[
\log p_i = z_i - \log\sum_j e^{z_j}.
\]

代入损失：

\[
\begin{aligned}
\ell &= -\sum_i y_i\Bigl(z_i - \log\sum_j e^{z_j}\Bigr) \\
&= -\sum_i y_i z_i + \Bigl(\sum_i y_i\Bigr)\log\sum_j e^{z_j} \\
&= -y^{\trans}z + \logsumexp(z).
\end{aligned}
\]

交叉熵裂成两项。第一项 \(-y^{\trans}z\)：logits 的线性函数，梯度直接是 \(-y\)。第二项 \(\logsumexp(z) = \log\sum_j e^{z_j}\)：它对标量 \(z_i\) 的偏导数是

\[
\frac{\partial}{\partial z_i}\log\sum_j e^{z_j}
= \frac{e^{z_i}}{\sum_j e^{z_j}}
= p_i,
\]

即 \(\nabla_z\logsumexp(z) = p\)。合起来：\(\nabla_z\ell = p - y\)。

两行推导。没有 Jacobian，没有 \(y_i/p_i\)，没有 \(K\times K\) 的中间量。关键在于 \textbf{log-softmax 这一步的合并}：把 softmax 的非线性与取对数的操作一步消化。交叉熵随后只是对这些 log-概率做线性组合——而 logits 中的线性项给出梯度 \(-y\)，LSE 的非线性项给出梯度 \(+p\)。

更一般地看：只有在交叉熵这个特定的损失函数下，\(\log p_i = z_i - \logsumexp(z)\) 才能把 \(p_i\) 中的指数从分母移出来，让损失对 logits 的依赖变成线性部分加 LSE 部分。如果把交叉熵换成均方误差 \(\|p-y\|^2\)，那条穿过 softmax Jacobian 的路径就避不开了——没有巧合可以消掉。交叉熵 + softmax 是目前分类任务的事实标准，不是没有原因的：除了概率解释和信息论背景外，它的梯度刚好是所有合理损失中计算最便宜的那一类。

\subsection{用一组具体数字验证一遍}

令 \(z=[2,1,0]\)，真实类别为第 0 类：\(y=[1,0,0]\)。

计算 softmax：\(\sum e^{z_j}=e^2+e^1+e^0\approx 7.389+2.718+1=11.107\)。
\[
p_0\approx 0.6653,\quad p_1\approx 0.2447,\quad p_2\approx 0.0900.
\]

交叉熵：\(\ell = -\log 0.6653\approx 0.4076\)。

梯度 \(\nabla_z\ell = p - y = [-0.3347,\;0.2447,\;0.0900]\)。

验证直觉：增大 \(z_0\)（正确类别的 logit），\(p_0\) 增大，损失减小——梯度 \(p_0-y_0=-0.3347\) 为负，反向传播会推高 \(z_0\)。增大 \(z_1\)（错误类别），\(p_1\) 增大，挤占 \(p_0\) 的空间，损失增大——梯度 \(p_1-y_1=0.2447\) 为正，反向传播会压低 \(z_1\)。

每个梯度分量只涉及该分量的预测和标签，不依赖其他分量的交叉项——这解释了为什么 \(O(K^2)\) 的 Jacobian 最终坍缩为 \(O(K)\) 的梯度：交叉项在求和时被 \(\sum_i y_i=1\) 这把梳子理干净了。

\subsection{Hessian 与平移不变性}

二阶导数更有意思：

\[
\nabla_z^2\ell = \diag(p) - pp^{\trans} \succeq 0.
\]

它是半正定的——交叉熵作为 logits 的函数是凸函数。对任意 \(v\in\R^K\)：

\[
v^{\trans}Hv = \sum_i p_i v_i^2 - \Bigl(\sum_i p_i v_i\Bigr)^2 = \Var_{i\sim p}(v_i) \ge 0.
\]

把 \(p_i\) 看作类别上的概率分布，二次型就是 \(v\) 在此分布下的方差。方差恒非负——这是 Hessian 半正定的一个漂亮概率论解释。

沿全一向量 \(\mathbf{1}\)：\(p\) 分布在常数方向上的方差为零。对应地，\(H\mathbf{1} = p - p(\sum_i p_i) = p-p = 0\)。同时梯度在这个方向也是零：\((p-y)^{\trans}\mathbf{1} = 1-1=0\)。这不是巧合：给所有 logits 同时加同一个常数 \(\Delta\)，softmax 输出不变——分子分母同时乘以 \(e^\Delta\) 相消。因此交叉熵对 logits 的常值平移不敏感，Hessian 的零空间正好是这个平移方向。

优化时这个平坦方向无伤大雅——梯度下降沿它不会动——但如果做 Newton 法或自然梯度下降，需要用伪逆或阻尼来处置这个零特征值。

\subsection{批量分类器：从单样本到矩阵}

设 \(X\in\R^{n\times d}\)，\(W\in\R^{d\times K}\)，\(b\in\R^K\)。logits：

\[
Z = XW + \mathbf{1}b^{\trans}\in\R^{n\times K}.
\]

令 \(P\in\R^{n\times K}\) 逐行 softmax，\(Y\) 为逐行标签。平均交叉熵损失对 logits 的梯度，由单样本结果逐行代入再除以 \(n\)：

\[
\bar Z = \frac1n(P-Y).
\]

用加横线的字母表示标量损失对某中间量的梯度，形状与该中间量一致。

全微分：\(dL = \trace(\bar Z^{\trans}dZ) = \trace\bigl(\bar Z^{\trans}(X\,dW + \mathbf{1}\,db^{\trans})\bigr)\)。

拆分两项。来自 \(dW\)：

\[
\trace(\bar Z^{\trans}X\,dW) = \trace\bigl((X^{\trans}\bar Z)^{\trans}dW\bigr)
\;\Longrightarrow\;
\nabla_W L = X^{\trans}\bar Z = \frac1n X^{\trans}(P-Y).
\]

形状：\(X^{\trans}\) 是 \(d\times n\)，\(\bar Z\) 是 \(n\times K\)，乘积 \(d\times K\)，与 \(W\) 一致。

来自 \(db\)：

\[
\trace(\bar Z^{\trans}\mathbf{1}\,db^{\trans})
= \trace(db^{\trans}\bar Z^{\trans}\mathbf{1})
= (\bar Z^{\trans}\mathbf{1})^{\trans}db
\;\Longrightarrow\;
\nabla_b L = \bar Z^{\trans}\mathbf{1} = \frac1n (P-Y)^{\trans}\mathbf{1}.
\]

形状为 \(K\)——沿样本轴的求和，把 \(n\times K\) 压缩到 \(K\) 维。

\subsection{数值稳定：减最大值是公式的血肉，不是后加的补丁}

直接算 \(e^{z_i}\)：\(z_i=1000\) 时 \(e^{1000}\) 超出 float32 可表示范围（\(\approx 3.4\times 10^{38}\)），溢出为 \(+\infty\)。恒等式救场：

\[
\logsumexp(z) = m + \log\sum_j e^{z_j-m},\qquad m=\max_j z_j.
\]

所有指数项 \(\le 0\)，最大项恰为 \(e^0=1\)，无上溢。log-softmax 应直接计算

\[
z_i - \logsumexp(z),
\]

永远不先算出 \(p_i\)（某些分量可能下溢为零）再取 \(\log\)。\(\log(0)=-\infty\)，而 \(z_i-\logsumexp(z)\) 不会碰这个问题——即使某个 \(z_i\) 远小于最大值，它的 log-概率只是一个绝对值较大的负数，所有运算安全地留在 log 域。

带有 mask 的注意力或序列模型中，无效位置的 logits 常被设为 \(-\infty\)。每行至少保留一个有效位置时 softmax 可正常归一化；若整行全被 mask，不存在合法概率分布——这不是 softmax 能处理的，必须由调用方在外部显式处理空 mask 行。

低精度训练（float16、bfloat16）再加一层微妙性：减最大值防了上溢，却无法恢复低精度下已被丢弃的微小差异。两个彼此很接近的 logits，在 float16 中可能变成相等的数，指数后相等，归一化就丢掉分辨力。在关键归约操作（max、sum）中内部使用 float32，是当前框架的标准妥协。

\subsection{标签平滑与温度：公式没变，问法变了}

标签平滑把 one-hot \(y\) 换成

\[
\tilde y = (1-\varepsilon)y + \varepsilon q,
\]

\(q\) 常取均匀分布 \(q_i=1/K\)。全部推导只用到 \(\sum_i y_i=1\)，\(\tilde y\) 依然满足，梯度仍是

\[
\nabla_z\ell = p - \tilde y.
\]

形式不变，但最优预测变了：不再要求正确类别的概率趋近 1，而是趋向 \((1-\varepsilon) + \varepsilon/K\)。这改变了模型的校准行为（预测概率不再过度自信），而没有改动反向传播公式本身。

温度 \(T>0\) 改变的是计算图的前向：

\[
p = \softmax(z/T).
\]

链式法则因此多出一个因子：

\[
\nabla_z\ell = \frac1T(p - y).
\]

高温 \(T\gg1\)：softmax 趋近均匀，\(p\approx 1/K\)，梯度幅度缩小——每步更新更谨慎。低温 \(T\ll1\)：softmax 趋近 argmax，梯度只在最大 logit 附近非零，其他位置近乎零——训练不稳定，方差大。忘记在梯度中补上 \(1/T\) 会让梯度尺度错位，这是实践中常见的坑：看上去只改了一行前向代码，反向也必须跟着变。任何``和标准公式看起来一样''的表达式，必须先对抗实际计算图中的每一个变换。
